\documentclass[12pt,a4paper]{article}

\usepackage[a4paper,margin=1in]{geometry}
\usepackage{graphicx}
\usepackage{booktabs}
\usepackage{amsmath,amssymb}
\usepackage{float}
\usepackage{hyperref}
\usepackage{xcolor}
\usepackage{listings}
\usepackage{enumitem}
\usepackage{fancyhdr}
\usepackage{longtable}
\usepackage{array}

\hypersetup{colorlinks=true,linkcolor=blue,urlcolor=blue}

\pagestyle{fancy}
\fancyhf{}
\lhead{ICS1512}
\rhead{Machine Learning Algorithms Laboratory}
\cfoot{\thepage}

\lstset{
language=Python,
basicstyle=\ttfamily\small,
numbers=left,
frame=single,
breaklines=true,
showstringspaces=false
}

\begin{document}

\begin{tabular}{|p{4cm}|p{11cm}|}
\hline
Experiment No. & 1 \\ \hline
Experiment Title & Working with Python Packages -- NumPy, Pandas, SciPy, Scikit-Learn, Matplotlib
\footnote{\textbf{GitHub Repository: }\url{https://github.com/username/repository}}\\ \hline
Student Name & Barath \\ \hline
Register Number & \\ \hline
Faculty & Dr. Poreddy Ajay Kumar Reddy \\ \hline
Submission Date & \\ \hline
\end{tabular}

\section{Objective}
To explore the core Python libraries used in a Machine Learning workflow -- NumPy, Pandas,
SciPy, Scikit-learn, and Matplotlib -- and to understand key operations such as array
manipulation, data preprocessing, mathematical computing, model building, and data
visualization. Further, to explore publicly available datasets from the UCI Machine Learning
Repository and Kaggle, identify the type of Machine Learning task associated with each
(Supervised/Unsupervised, Classification/Regression), and perform Exploratory Data Analysis
(EDA) on them using a reusable, general-purpose EDA toolkit.

\section{Problem Statement}
Given a set of five publicly available datasets -- the Iris dataset, a Loan Amount Prediction
dataset, a Diabetes Prediction dataset, an Email Spam Classification dataset, and the MNIST
Handwritten Digit dataset -- the task is to:
\begin{itemize}
\item Identify the nature of each dataset (tabular numeric, tabular categorical, text, or
image) and the corresponding Machine Learning task (regression, binary classification,
multi-class classification).
\item Load each dataset using Pandas and inspect it using NumPy- and SciPy-based numerical
routines.
\item Build a single, reusable EDA toolkit capable of handling all four data modalities
(tabular, text, image-as-CSV, and image-folder datasets) so that the same code can be applied
to any new dataset with minimal changes.
\item Summarize, visualize, and interpret each dataset so that an appropriate feature
selection technique and ML algorithm can be recommended for it.
\end{itemize}
The input to the workflow is a raw dataset file (CSV or an image directory); the output is a
statistical and visual understanding of the dataset (distributions, missing values,
correlations, class balance, outliers) that guides subsequent preprocessing and modelling
decisions.

\section{Python Libraries for Machine Learning}

\subsection{NumPy}
NumPy provides the \texttt{ndarray} object for fast, vectorized numerical computation.
Functions explored include \texttt{np.array}, \texttt{np.reshape}, \texttt{np.arange},
\texttt{np.linspace}, \texttt{np.mean}, \texttt{np.std}, \texttt{np.sum}, broadcasting rules,
boolean indexing (used for IQR-based outlier detection), and \texttt{np.number} for
selecting numeric dtypes from a DataFrame.

\subsection{Pandas}
Pandas provides the \texttt{DataFrame}/\texttt{Series} structures used to load and manipulate
tabular data. Key operations explored: \texttt{read\_csv}, \texttt{head}, \texttt{shape},
\texttt{dtypes}, \texttt{info}, \texttt{describe}, \texttt{isnull().sum()},
\texttt{duplicated()}, \texttt{nunique()}, \texttt{value\_counts()},
\texttt{select\_dtypes()}, \texttt{corr()}, \texttt{skew()}, and \texttt{kurt()} --
all of which are used directly inside the EDA toolkit in Section~\ref{sec:eda-toolkit}.

\subsection{SciPy}
SciPy extends NumPy with statistical and scientific routines. Relevant modules for this
experiment include \texttt{scipy.stats} (e.g. \texttt{skew}, \texttt{kurtosis},
\texttt{zscore}, \texttt{chi2\_contingency} for categorical association tests, and
\texttt{pearsonr}/\texttt{spearmanr} for correlation testing), which complement Pandas'
built-in statistical summaries.

\subsection{Scikit-learn}
Scikit-learn provides a consistent \texttt{fit}/\texttt{transform}/\texttt{predict} API for
preprocessing, feature selection, and modelling. Relevant utilities include
\texttt{train\_test\_split}, \texttt{StandardScaler}/\texttt{MinMaxScaler},
\texttt{LabelEncoder}/\texttt{OneHotEncoder}, \texttt{SelectKBest}, \texttt{chi2}, \texttt{f\_classif}
(ANOVA F-test), and estimator classes such as \texttt{LogisticRegression},
\texttt{DecisionTreeClassifier}, \texttt{RandomForestClassifier}, and \texttt{KMeans}.

\subsection{Matplotlib \& Seaborn}
Matplotlib (via \texttt{pyplot}) and Seaborn provide the visualization layer used throughout
the toolkit -- histograms (\texttt{sns.histplot}), box plots (\texttt{sns.boxplot}), count
plots (\texttt{sns.countplot}), heatmaps (\texttt{sns.heatmap}) for missing values and
correlation, and \texttt{sns.pairplot} for multivariate numeric relationships.

\section{Universal EDA Toolkit}
\label{sec:eda-toolkit}
To avoid rewriting EDA code for every new dataset, a single reusable toolkit was built with
four entry-point functions, one for each data modality encountered in this experiment:
\texttt{eda\_tabular()} for numeric/categorical CSV data (Loan, Diabetes, Iris),
\texttt{eda\_text()} for text CSV data (Email Spam), \texttt{eda\_image\_csv()} for
pixel-flattened image CSV data (MNIST), and \texttt{eda\_image\_folder()} for image datasets
stored as class-labelled folders (Handwritten Character Recognition).

\subsection{Setup}
\begin{lstlisting}
import pandas as pd
import numpy as np
import matplotlib.pyplot as plt
import seaborn as sns
import os
from PIL import Image
\end{lstlisting}

\subsection{Tabular EDA -- \texttt{eda\_tabular()}}
Used for datasets such as \textit{Loan Amount Prediction}, \textit{Diabetes Prediction}, and
the \textit{Iris} dataset. It prints the head, shape, column names, data types, a full
\texttt{describe()} summary, missing-value counts, and duplicate-row counts; plots a missing
value heatmap, histograms and box plots for every numeric column, and count plots for every
low-cardinality categorical column; computes a correlation heatmap and pair plot when there
are multiple numeric columns; and finally reports skewness, kurtosis, and IQR-based outlier
counts for every numeric feature.
\begin{lstlisting}
def eda_tabular(csv_path):
    df = pd.read_csv(csv_path)

    print("="*60)
    print("HEAD")
    print(df.head())

    print("\nShape:", df.shape)
    print("\nColumns:", list(df.columns))
    print("\nData Types")
    print(df.dtypes)

    print("\nInfo")
    df.info()

    print("\nSummary")
    print(df.describe(include="all"))

    print("\nMissing Values")
    print(df.isnull().sum())

    sns.heatmap(df.isnull(), cbar=False)
    plt.title("Missing Values")
    plt.show()

    print("\nDuplicates:", df.duplicated().sum())

    # ----------------------------- # Unique Values # -----------------------------
    print("\n"+"="*60)
    print("UNIQUE VALUES")
    print("="*60)
    for col in df.columns: print(f"{col} : {df[col].nunique()}")

    num_cols = df.select_dtypes(include=np.number).columns
    cat_cols = df.select_dtypes(exclude=np.number).columns

    for c in num_cols:
        plt.figure(figsize=(6,4))
        sns.histplot(df[c].dropna(), kde=True)
        plt.title(c)
        plt.show()

        plt.figure(figsize=(6,2.5))
        sns.boxplot(x=df[c].dropna())
        plt.title(c)
        plt.show()

    for c in cat_cols:
        if df[c].nunique()<=20:
            plt.figure(figsize=(7,4))
            sns.countplot(x=df[c])
            plt.xticks(rotation=45)
            plt.title(c)
            plt.show()

    if len(num_cols)>1:
        plt.figure(figsize=(10,8))
        sns.heatmap(df[num_cols].corr(),annot=True,cmap="coolwarm")
        plt.show()

        if len(num_cols)<=8:
            sns.pairplot(df[num_cols])
            plt.show()

    print("\nSkewness")
    print(df[num_cols].skew())

    print("\nKurtosis")
    print(df[num_cols].kurt())

    print("\nOutliers (IQR)")
    for c in num_cols:
        q1=df[c].quantile(.25)
        q3=df[c].quantile(.75)
        iqr=q3-q1
        low=q1-1.5*iqr
        high=q3+1.5*iqr
        print(c,((df[c]<low)|(df[c]>high)).sum())
\end{lstlisting}

\subsection{Text EDA -- \texttt{eda\_text()}}
Used for the \textit{Email Spam} dataset. It reports the standard shape/type/missing-value
summary, then treats low-cardinality object columns (e.g. the spam/ham label) as categorical
and plots their distribution, while high-cardinality object columns (the raw message text)
are analyzed for character length and word-count distributions.
\begin{lstlisting}
def eda_text(csv_path):
    df=pd.read_csv(csv_path)

    print(df.head())
    print(df.shape)
    print(df.dtypes)
    print(df.isnull().sum())
    print("Duplicates:",df.duplicated().sum())

    text_cols=df.select_dtypes(include="object").columns

    print("Text Columns:",list(text_cols))

    for c in text_cols:
        if df[c].nunique()<=20:
            print(df[c].value_counts())
            sns.countplot(x=df[c])
            plt.xticks(rotation=45)
            plt.title(c)
            plt.show()

    for c in text_cols:
        if df[c].nunique()>20:
            temp=df[c].astype(str)
            lengths=temp.apply(len)
            words=temp.apply(lambda x:len(x.split()))

            plt.figure(figsize=(6,4))
            sns.histplot(lengths,bins=30,kde=True)
            plt.title("Character Length")
            plt.show()

            plt.figure(figsize=(6,4))
            sns.histplot(words,bins=30,kde=True)
            plt.title("Word Count")
            plt.show()

            print(lengths.describe())
            print(words.describe())
            break
\end{lstlisting}

\subsection{Image EDA -- \texttt{eda\_image\_csv()} and \texttt{eda\_image\_folder()}}
\texttt{eda\_image\_csv()} is used for the \textit{MNIST} dataset, where each row is a
flattened image with the label in the first column: it plots the class distribution, renders
a $3\times3$ grid of sample digit images, and reports pixel-intensity mean, standard deviation,
and distribution. \texttt{eda\_image\_folder()} is used for the \textit{Handwritten Character
Recognition} dataset when it is organized as one sub-folder per class: it counts images per
class and displays one sample image from each of the first six classes.
\begin{lstlisting}
def eda_image_csv(csv_path):
    df=pd.read_csv(csv_path)

    print(df.head())
    print(df.shape)
    print(df.isnull().sum())

    target=df.columns[0]

    sns.countplot(x=df[target])
    plt.title("Class Distribution")
    plt.show()

    size=int((df.shape[1]-1)**0.5)

    plt.figure(figsize=(8,8))
    for i in range(min(9,len(df))):
        plt.subplot(3,3,i+1)
        img=df.iloc[i,1:].values.reshape(size,size)
        plt.imshow(img,cmap="gray")
        plt.title(df.iloc[i,0])
        plt.axis("off")
    plt.tight_layout()
    plt.show()

    pixels=df.iloc[:,1:]
    print("Pixel Mean:",pixels.mean().mean())
    print("Pixel Std:",pixels.std().mean())

    sns.histplot(pixels.values.flatten(),bins=50)
    plt.title("Pixel Distribution")
    plt.show()


def eda_image_folder(folder_path):
    classes=[c for c in os.listdir(folder_path) if os.path.isdir(os.path.join(folder_path,c))]
    print("Classes:",classes)

    for c in classes:
        print(c,":",len(os.listdir(os.path.join(folder_path,c))),"images")

    plt.figure(figsize=(12,6))
    for i,c in enumerate(classes[:6]):
        folder=os.path.join(folder_path,c)
        img=Image.open(os.path.join(folder,os.listdir(folder)[0]))
        plt.subplot(2,3,i+1)
        plt.imshow(img)
        plt.title(c)
        plt.axis("off")
    plt.tight_layout()
    plt.show()
\end{lstlisting}

\subsection{Example Usage}
\begin{lstlisting}
# Tabular
eda_tabular("loan.csv")

# Text
eda_text("spam.csv")

# Image CSV (MNIST)
eda_image_csv("mnist.csv")

# Image Folder
eda_image_folder("dataset/train")
\end{lstlisting}

\textbf{Inference:} A single, modality-aware toolkit removes the need to hand-write EDA code
for every new dataset. Because each function follows the same pattern -- load, summarize,
visualize, quantify -- it standardizes the analysis across very different data types
(structured numeric/categorical data, free text, and images), making it easy to plug in any
of the five datasets used in this experiment by simply calling the matching function.

\section{Dataset Description}
The five datasets below were obtained as follows: Iris via \texttt{sklearn.datasets.load\_iris()};
the Pima Indians Diabetes, Loan Prediction, and SMS Spam Collection datasets from their
respective public GitHub mirrors of the original UCI/Kaggle releases; and MNIST from a public
CSV mirror of the standard handwritten-digit dataset (a 3,000-row sample used here for
tractable EDA). All five were loaded and profiled using the toolkit in Section~\ref{sec:eda-toolkit}.

\begin{longtable}{|p{3.6cm}|p{2.6cm}|p{2.6cm}|p{2.5cm}|p{2.5cm}|}
\hline
\textbf{Dataset} & \textbf{Samples $\times$ Features} & \textbf{Missing Values} & \textbf{Duplicates} & \textbf{Classes} \\ \hline
\endhead
Iris & 150 $\times$ 5 & 0 & 1 & 3 (balanced, 50 each) \\ \hline
Pima Diabetes & 768 $\times$ 9 & 0 (nulls); 5 zero-Glucose, 374 zero-Insulin & 0 & 2 (500 / 268) \\ \hline
Loan Prediction & 614 $\times$ 13 & 149 cells across 7 columns (worst: Credit\_History, 50) & 0 & 2 (422 Y / 192 N) \\ \hline
SMS Spam & 5,572 $\times$ 2 & 0 & 403 & 2 (4,825 ham / 747 spam) \\ \hline
MNIST (sample) & 3,000 $\times$ 785 & 0 & -- & 10 (roughly balanced, 270--332 each) \\ \hline
\end{longtable}

\section{Dataset Exploration and ML Task Identification}

\begin{longtable}{|p{3.3cm}|p{2.6cm}|p{3.6cm}|p{4.7cm}|}
\hline
\textbf{Dataset} & \textbf{Type of ML Task} & \textbf{Feature Selection Technique} & \textbf{Suitable ML Algorithm} \\ \hline
\endhead
Iris Dataset & Supervised -- Multi-class Classification & ANOVA F-test (\texttt{SelectKBest} with \texttt{f\_classif}) & Logistic Regression / K-Nearest Neighbors / Decision Tree \\ \hline
Loan Amount Prediction & Supervised -- Regression & Correlation analysis, ANOVA F-test (\texttt{f\_regression}) & Linear Regression / Random Forest Regressor \\ \hline
Predicting Diabetes & Supervised -- Binary Classification & Chi-square test (\texttt{SelectKBest} with \texttt{chi2}) & Logistic Regression / Random Forest Classifier \\ \hline
Classification of Email Spam & Supervised -- Binary Classification (Text) & Chi-square test on TF-IDF/Bag-of-Words features & Naive Bayes / Support Vector Machine (SVM) \\ \hline
Handwritten Character Recognition / MNIST & Supervised -- Multi-class Classification (Image) & Variance thresholding / PCA (dimensionality reduction over pixel features) & Convolutional Neural Network (CNN) / SVM / Random Forest \\ \hline
\end{longtable}

\textbf{Inference:} All five datasets correspond to supervised learning tasks, split between
regression (Loan Amount) and classification (Iris, Diabetes, Spam, MNIST). The choice of
feature selection technique depends on the feature type: chi-square and ANOVA suit
categorical/numeric tabular features, TF-IDF-based statistics suit text, and
variance-based/PCA techniques suit high-dimensional pixel data.

\section{Exploratory Data Analysis -- Visualizations}
The toolkit in Section~\ref{sec:eda-toolkit} was run on all five datasets described above.
The resulting figures and the statistics they produced are reported per dataset below.

\subsection{Iris Dataset}

\begin{figure}[H]
\centering
\includegraphics[width=0.55\linewidth]{iris_species_count.png}
\caption{Species count plot -- Iris}
\end{figure}
\textbf{Inference:} The three species (setosa, versicolor, virginica) are perfectly balanced
at 50 samples each. This balance means accuracy is a reliable evaluation metric here and no
class-imbalance handling (e.g. resampling) is required before training a classifier.

\begin{figure}[H]
\centering
\includegraphics[width=0.6\linewidth]{iris_corr.png}
\caption{Correlation heatmap -- Iris}
\end{figure}
\textbf{Inference:} Petal length and petal width are very strongly correlated with each
other, and both correlate strongly with sepal length, while sepal width correlates weakly
and even negatively with the other features. This indicates that petal measurements alone
carry most of the discriminative signal for species classification, and the two petal
features are partially redundant with one another.

\begin{figure}[H]
\centering
\includegraphics[width=0.85\linewidth]{iris_pairplot.png}
\caption{Pair plot coloured by species -- Iris}
\end{figure}
\textbf{Inference:} Setosa is linearly separable from the other two species on almost every
feature pair, especially petal length vs. petal width, while versicolor and virginica show
some overlap. This suggests a simple linear or distance-based classifier (Logistic
Regression, KNN) should already perform well, with most residual error expected between
versicolor and virginica.

\begin{figure}[H]
\centering
\includegraphics[width=0.45\linewidth]{iris_box_sepal_width.png}
\caption{Sepal width box plot -- Iris}
\end{figure}
\textbf{Inference:} The IQR method flags 4 outlier values in sepal width out of 150 samples
(about 2.7\%). Given the small number and that Iris is a well-curated benchmark dataset,
these are most plausibly genuine measurement variation rather than data-entry errors, so
retaining them is reasonable.

\subsection{Diabetes Dataset}

\begin{figure}[H]
\centering
\includegraphics[width=0.55\linewidth]{diabetes_hist_glucose.png}
\caption{Glucose distribution -- Diabetes}
\end{figure}
\textbf{Inference:} Glucose is close to a bell-shaped distribution but with a noticeable
spike of zero values -- 5 records have Glucose $=0$, which is physiologically impossible and
indicates the value is actually a missing-data placeholder rather than a true zero reading.

\begin{figure}[H]
\centering
\includegraphics[width=0.55\linewidth]{diabetes_class_count.png}
\caption{Outcome class distribution -- Diabetes}
\end{figure}
\textbf{Inference:} Outcome is imbalanced (500 non-diabetic vs. 268 diabetic, roughly 65:35).
This is mild but non-trivial imbalance, so precision, recall, and F1-score for the positive
class should be reported alongside accuracy, and class weighting or resampling could improve
recall on the minority (diabetic) class.

\begin{figure}[H]
\centering
\includegraphics[width=0.45\linewidth]{diabetes_box_insulin.png}
\caption{Insulin box plot -- Diabetes}
\end{figure}
\textbf{Inference:} The IQR test flags 34 outliers in Insulin. Combined with the fact that
374 of 768 records (about 49\%) have Insulin exactly equal to 0, this again points to
placeholder zeros for unrecorded measurements rather than genuine physiological readings --
this column will need explicit missing-value imputation rather than outlier removal.

\begin{figure}[H]
\centering
\includegraphics[width=0.75\linewidth]{diabetes_corr.png}
\caption{Correlation heatmap -- Diabetes}
\end{figure}
\textbf{Inference:} Glucose shows the strongest correlation with Outcome (Pearson $r \approx
0.47$) of any single feature, consistent with its clinical role in diabetes diagnosis. Other
features (BMI, Age, Pregnancies) show weaker but still positive associations, while no pair
of predictors is correlated highly enough to suggest serious multicollinearity.

\subsection{Loan Prediction Dataset}

\begin{figure}[H]
\centering
\includegraphics[width=0.7\linewidth]{loan_missing.png}
\caption{Missing value heatmap -- Loan Prediction}
\end{figure}
\textbf{Inference:} Missing values are spread across seven columns, most severely
Credit\_History (50 missing), Self\_Employed (32), and LoanAmount (22). Credit\_History is
typically the single strongest predictor of loan approval, so its missingness needs careful
handling (mode/most-frequent imputation, or a separate ``unknown'' category) rather than
row deletion, which would discard a meaningful fraction of the 614 records.

\begin{figure}[H]
\centering
\includegraphics[width=0.55\linewidth]{loan_hist_income.png}
\caption{Applicant income distribution -- Loan Prediction}
\end{figure}
\textbf{Inference:} Applicant income is heavily right-skewed (skewness $\approx 6.5$), with
most applicants clustered at lower incomes and a long tail of high earners. This strongly
suggests a log transformation before feeding income into a regression model, and explains
why 50 values are flagged as IQR outliers on the raw scale.

\begin{figure}[H]
\centering
\includegraphics[width=0.55\linewidth]{loan_status_count.png}
\caption{Loan status distribution -- Loan Prediction}
\end{figure}
\textbf{Inference:} Loan\_Status is imbalanced, with about 69\% of applications approved (Y:
422) versus 31\% rejected (N: 192). This imbalance should be reflected in the evaluation
metric choice (F1/recall on the rejected class matter as much as overall accuracy) if the
downstream task is framed as approval classification rather than amount regression.

\begin{figure}[H]
\centering
\includegraphics[width=0.6\linewidth]{loan_property_area.png}
\caption{Property area distribution -- Loan Prediction}
\end{figure}
\textbf{Inference:} Applications are reasonably distributed across the three property-area
categories (Urban, Semiurban, Rural), with no single category dominating -- so one-hot
encoding this feature should not introduce severe sparsity or imbalance concerns.

\subsection{Email Spam Dataset}

\begin{figure}[H]
\centering
\includegraphics[width=0.5\linewidth]{spam_class_count.png}
\caption{Ham vs. spam class distribution}
\end{figure}
\textbf{Inference:} The dataset is strongly imbalanced -- 4,825 ham messages versus only 747
spam (about 87:13). A classifier trained naively on raw accuracy could reach $\sim$87\% by
always predicting ``ham,'' so precision/recall/F1 on the spam class, or techniques such as
class weighting, are essential for a meaningful evaluation.

\begin{figure}[H]
\centering
\includegraphics[width=0.55\linewidth]{spam_char_len.png}
\caption{Character length distribution -- Spam}
\end{figure}

\begin{figure}[H]
\centering
\includegraphics[width=0.55\linewidth]{spam_word_count.png}
\caption{Word count distribution -- Spam}
\end{figure}
\textbf{Inference:} Message length is itself discriminative: the toolkit's grouped statistics
show spam messages average about 139 characters versus roughly 71 for ham -- almost double.
Both distributions are right-skewed with a long tail toward longer messages (max 910
characters), so character/word count could be added as an engineered feature alongside
TF-IDF text features.

\subsection{MNIST Dataset (3,000-sample subset)}

\begin{figure}[H]
\centering
\includegraphics[width=0.55\linewidth]{mnist_class_count.png}
\caption{Digit class distribution -- MNIST}
\end{figure}
\textbf{Inference:} The 10 digit classes are close to balanced in this sample, ranging from
270 (digit 5) to 332 (digit 1) examples -- roughly a 20\% spread. This near-balance means
accuracy remains a reasonable primary metric, though a per-class confusion matrix is still
useful given digits like 4/9 or 3/8 are visually similar and prone to confusion.

\begin{figure}[H]
\centering
\includegraphics[width=0.7\linewidth]{mnist_samples.png}
\caption{Sample digit images -- MNIST}
\end{figure}
\textbf{Inference:} The sample grid confirms the images are correctly labelled, centered,
grayscale $28\times28$ handwritten digits consistent with the standard MNIST format, with
visible natural variation in stroke width and slant across samples of the same digit.

\begin{figure}[H]
\centering
\includegraphics[width=0.55\linewidth]{mnist_pixel_dist.png}
\caption{Pixel intensity distribution -- MNIST}
\end{figure}
\textbf{Inference:} Pixel values range from 0 to 255 with a mean of about 33 and a large
standard deviation ($\approx$78), and only about 19\% of pixels are non-zero. This confirms
the images are mostly dark background with sparse bright strokes, and supports normalizing
pixel values (dividing by 255) before feeding them into a CNN or any distance-based model.

\section{Machine Learning Workflow}
The general workflow followed for each dataset consists of:
\begin{enumerate}
\item \textbf{Loading the dataset} -- reading the raw file with \texttt{pandas.read\_csv()}
(or \texttt{os}/\texttt{PIL} for image folders).
\item \textbf{Exploratory Data Analysis and Visualization} -- using the toolkit in
Section~\ref{sec:eda-toolkit} to inspect summary statistics, histograms, bar charts, box
plots, and correlation heatmaps.
\item \textbf{Data Preprocessing} -- handling missing values, dropping irrelevant features,
encoding categorical variables (\texttt{LabelEncoder}/\texttt{OneHotEncoder}), and
normalizing/standardizing numeric features (\texttt{StandardScaler}/\texttt{MinMaxScaler}).
\item \textbf{Feature Selection} -- applying \texttt{SelectKBest} with the chi-square test or
ANOVA F-test to retain the most relevant features, as summarized in the table above.
\item \textbf{Data Splitting} -- dividing each dataset into training, validation, and testing
subsets using \texttt{train\_test\_split}.
\item \textbf{Performance Evaluation} -- selecting metrics appropriate to the task (accuracy,
precision, recall, F1-score, ROC-AUC for classification; RMSE/MAE/$R^2$ for regression) and
visualizing results with confusion matrices or residual plots.
\end{enumerate}

\section{Experimental Setup}
\begin{tabular}{ll}
\toprule
Python Version & \\
Pandas / NumPy / SciPy / Scikit-Learn Version & \\
IDE & Jupyter Notebook \\
Operating System & \\
Hardware & \\
\bottomrule
\end{tabular}

\section{Discussion}
The Universal EDA Toolkit demonstrates that a small set of well-structured functions can
generalize across tabular, text, and image data, which are the three broad modalities
encountered across the five prescribed datasets. Grouping missing-value analysis, univariate
distribution plots, and correlation/outlier analysis into a single tabular routine
significantly reduces repeated exploration code. A limitation is that the image-CSV function
assumes square, single-channel images and that the text function only profiles the first
high-cardinality text column found -- both would need to be generalized further for more
complex, multi-column or multi-channel datasets.

\section{Conclusion}
This experiment explored the core NumPy, Pandas, SciPy, Scikit-learn, and Matplotlib/Seaborn
functionality needed for a Machine Learning workflow, and used it to identify the ML task
type, suitable feature selection technique, and suitable algorithm for five public datasets
(Iris, Loan Amount Prediction, Diabetes Prediction, Email Spam Classification, and MNIST /
Handwritten Character Recognition). A reusable, modality-aware EDA toolkit was built and
applied to standardize the exploratory analysis across all five datasets.

\section{References}
\begin{enumerate}[label={[\arabic*]}]
\item NumPy Documentation, \url{https://numpy.org/doc/}
\item Pandas Documentation, \url{https://pandas.pydata.org/docs/}
\item SciPy Documentation, \url{https://docs.scipy.org/doc/scipy/}
\item Scikit-learn Documentation, \url{https://scikit-learn.org/stable/}
\item Matplotlib Documentation, \url{https://matplotlib.org/stable/}
\item UCI Machine Learning Repository, \url{https://archive.ics.uci.edu/}
\item Kaggle Datasets, \url{https://www.kaggle.com/datasets}
\end{enumerate}

\end{document}
